\mwnewpart
\renewcommand\runningtitle{Section 1. Introduction}
\chapter{Introduction}
\sectioncontents

\section{What RetrOS/LT Is}\ixs{RetrOS/LT}

RetrOS/LT is a disk operating system for the IBM Personal Computer, the
IBM Personal Computer XT, and compatible computers. It starts the
computer, keeps your files on diskettes and fixed disks, and runs your
programs. ``LT'' stands for \emph{lite}: RetrOS/LT is the small member of
the RetrOS family, built to be what the IBM Personal Computer ran in 1983,
and nothing more.

RetrOS/LT is derived from the MS-DOS 2 source code that Microsoft
Corporation released under the MIT License. It is made to behave as IBM
PC DOS 2.10 does on the same computers: the same commands, with the same
switches, the same messages and the same results. Programs written for DOS
2 run under it unchanged. To a program that asks, RetrOS/LT answers that
it is DOS version 2.10.

When the computer starts from a diskette with no \file{AUTOEXEC.BAT}
file, RetrOS/LT asks the date and time, shows its name, and then the
prompt:

\begin{dialogue}
Current date is Tue  1-01-1980
Enter new date: 11-13-2026
Current time is  0:00:12.34
Enter new time: 9:30


RetrOS/LT Version 1.0
(C) Copyright Republic Lyon Mountain Corporation 2026
(C) Copyright Microsoft Corp. 1981, 1982, 1983

A>
\end{dialogue}

Type the year in four digits: RetrOS/LT, like DOS 2, takes a two-digit
year as 19\var{yy}, and accepts dates from 1980 through 2099.

\section{The Computers It Runs On}\ixs{computers}\ixs{IBM Personal Computer}

\begin{tabularx}{\textwidth}{@{}P{1.25in}>{\raggedright\arraybackslash}X@{}}
\toprule
Processor & 8088 or 8086\\
Computers & IBM Personal Computer, Personal Computer XT, Portable Personal
Computer, and compatibles\\
Memory & 64 KB or more. RetrOS/LT has been tested with 256 KB and 640 KB.\\
Diskette drives & One or two 5-1/4 inch drives, single- or double-sided\\
Fixed disk & Optional: one or two, through the IBM fixed disk adapter or one
that follows it\\
Display & Monochrome or Color/Graphics Monitor Adapter\\
\bottomrule
\end{tabularx}

\subsection{The IBM Personal Computer} To start from a fixed disk, the
Personal Computer needs its second BIOS, of October 1982, which looks for
the program on a fixed disk adapter. With it, the Personal Computer can
start from a fixed disk whose adapter is in the system unit, or in the IBM
Expansion Unit. A Personal Computer with the first BIOS, of April 1981,
runs RetrOS/LT from diskettes, and can use a fixed disk only after
starting from a diskette.

\subsection{One diskette drive} On a computer with one diskette drive, the
drive is both A and B. When a program asks for the drive it is not using
now, RetrOS/LT asks you to change diskettes:

\begin{dialogue}
Insert diskette for drive B: and strike
any key when ready
\end{dialogue}

\section{What Is on the Diskettes}\ixs{diskettes, RetrOS/LT}

RetrOS/LT comes on two 180 KB single-sided diskettes, which any drive of
the IBM Personal Computer or XT can read, as PC DOS 2.10 did:

\begin{itemize}
\item \textbf{The DOS diskette} starts the computer. It holds the system,
  the command processor \file{COMMAND.COM}, the commands and device
  drivers, and \typed{INSTALL}, which puts RetrOS/LT on a fixed disk. When
  you start the computer from it, it asks the date and time and then runs
  Install; press \key{Esc} to leave Install and use the commands from
  drive A.
\item \textbf{The Supplemental Programs diskette} holds the programmer's
  tools: \typed{DEBUG}, \typed{LINK} and \typed{EXE2BIN}.
\end{itemize}

The DOS diskette also holds \file{LICENSE.TXT}, the RetrOS/LT License,
and \file{NOTICE.TXT}, the notice that accompanies Microsoft's code. Both
are printed in Appendix~D. Appendix~A lists every file on both diskettes.

\caution{Make copies of the RetrOS/LT diskettes before you use them, and
keep the originals in a safe place. See ``Copying a Diskette'' in
Section~4.}

\section{How RetrOS/LT Differs from PC DOS 2.10}\ixs{PC DOS 2.10, differences from}

RetrOS/LT has the commands of PC DOS 2.10 and behaves as they do. The
differences are few, and you can see them:

\begin{itemize}
\item \textbf{The names.} RetrOS/LT shows its own name, and its two system
  files are \file{RETBIO.COM} and \file{RETDOS.COM} where PC DOS has
  \file{IBMBIO.COM} and \file{IBMDOS.COM}. For that reason the
  \typed{FORMAT} and \typed{SYS} of each system put only their own system
  files on a disk.
\item \textbf{Install.} PC DOS 2.10 had no installation program; RetrOS/LT
  adds one. See Section~2.
\item \textbf{FC.} RetrOS/LT includes \typed{FC}, Microsoft's file
  comparison program, beside IBM's \typed{COMP}.
\item \textbf{Starting up.} RetrOS/LT reads its system files a track at a
  time, so it starts a little faster from a diskette.
\end{itemize}

\section{What RetrOS/LT Does Not Include}

\begin{itemize}
\item \textbf{BASIC.} IBM's \typed{BASIC} and \typed{BASICA} build on the
  Cassette BASIC in the Personal Computer's read-only memory and are IBM's
  own programs; RetrOS/LT does not include them, nor IBM's sample BASIC
  programs.
\item \textbf{The features of later versions of DOS}, such as volume
  serial numbers, \typed{ATTRIB}, \typed{LABEL}, \typed{XCOPY} or
  international support. A program that needs DOS 3 or later says so and
  stops; see ``Programs That Check the DOS Version'' in Section~7.
\end{itemize}
